% ============================================================
% WORK_ONE: DTMF Dialer
% ============================================================

\subsection{WORK\_ONE: Sintetizador y verificación de tonos DTMF}

\subsubsection{Objetivo y metodología}

La tarea consistió en implementar un generador de señales DTMF
(\emph{Dual-Tone Multi-Frequency}) — el esquema de tonos usado por los
teclados telefónicos — sintetizando cada dígito como la suma de dos
senoidales (grupo de baja frecuencia, 697--941~Hz, y grupo de alta
frecuencia, 1209--1477~Hz) según el estándar ITU-T~Q.23. La implementación
(\texttt{dtmf\_dialer/dtmf.py}) mezcla ambas senoidales como su promedio
($0.5\,(\sin_1+\sin_2)$) para evitar recorte antes de cuantizar a PCM de
16~bits, con muestreo a 44.1~kHz. El paquete expone tanto una interfaz de
línea de comandos (\texttt{main.py}) como una GUI de escritorio en PyQt6
(\texttt{run\_gui.py}) que comparten el mismo backend de síntesis, y
produce archivos \texttt{.wav} reproducibles.

Se verificó la implementación en dos niveles: (1) la suite de pruebas
unitarias del proyecto (\texttt{tests/test\_dtmf.py}), y (2) un análisis
espectral independiente por FFT sobre uno de los archivos \texttt{.wav}
generados, para confirmar que las frecuencias realmente sintetizadas
coinciden con la tabla ITU-T.

\subsubsection{Resultados}

La suite de 15 pruebas unitarias (\texttt{pytest}) se ejecutó completa y
pasó en su totalidad, cubriendo la tabla de frecuencias, el conteo de
muestras y tipo de dato (\texttt{int16}) de los tonos generados, los
límites de amplitud, el escalado lineal de secuencias multi-dígito, el
descarte de caracteres no numéricos y la integridad del ciclo
escritura/lectura de \texttt{.wav}.

Para verificar el contenido espectral real, se tomó el archivo generado
\texttt{output/dial\_7004191.wav} (marcación del número \texttt{7004191})
y se aplicó FFT a cada bloque de 0.5~s correspondiente a un dígito,
buscando el pico dominante dentro de las bandas de baja (650--1000~Hz) y
alta frecuencia (1150--1500~Hz). Los picos detectados coincidieron con la
tabla ITU-T con una desviación máxima de 1~Hz, consistente con la
resolución espectral del bin de FFT usado ($\Delta f \approx 2$~Hz para
una ventana de 0.5~s a 44.1~kHz):

\begin{table}[!t]
\centering
\caption{Frecuencias detectadas por FFT vs. tabla ITU-T (\texttt{dial\_7004191.wav})}
\label{tab:dtmf_verif}
\small
\begin{tabular}{@{}c c c c c@{}}
\toprule
Dígito & $f_{low}$ detectada & $f_{high}$ detectada & Nominal ITU-T \\
\midrule
7 & 852.0 Hz & 1208.0 Hz & (852, 1209) \\
0 & 940.0 Hz & 1336.0 Hz & (941, 1336) \\
0 & 940.0 Hz & 1336.0 Hz & (941, 1336) \\
4 & 770.0 Hz & 1208.0 Hz & (770, 1209) \\
1 & 696.0 Hz & 1210.0 Hz & (697, 1209) \\
9 & 852.0 Hz & 1476.0 Hz & (852, 1477) \\
1 & 696.0 Hz & 1210.0 Hz & (697, 1209) \\
\bottomrule
\end{tabular}
\end{table}

La Fig.~\ref{fig:dtmf_spectrum} muestra el espectro de magnitud de cada uno
de los siete dígitos de la secuencia \texttt{7004191}, confirmando
visualmente que cada tono ocupa exactamente dos picos angostos en las
frecuencias esperadas, sin energía espuria significativa fuera de las
bandas DTMF.

\begin{figure}[!t]
\centering
\includegraphics[width=\columnwidth]{work_one_dtmf_spectrum.png}
\caption{Espectro FFT de cada dígito de la secuencia \texttt{7004191} generada por el sintetizador DTMF, con los dos picos de frecuencia (grupo bajo y alto) claramente resueltos.}
\label{fig:dtmf_spectrum}
\end{figure}

\subsubsection{Limitaciones}

El proyecto no incluye un módulo de generación de figuras propio; la
Fig.~\ref{fig:dtmf_spectrum} fue producida específicamente para este
informe mediante un análisis FFT ad-hoc sobre un archivo \texttt{.wav}
real ya generado por el propio programa, no mediante datos sintéticos
inventados. La reproducción de audio (\texttt{play\_wav}, basada en
\texttt{winsound}) es específica de Windows y, por diseño del proyecto,
está excluida de las pruebas automatizadas; solo se validó la síntesis y
exportación, que son multiplataforma.
